\section{モデリングの目から見た検定3；多群の平均値差を求めるモデル}
\subsection{授業内容}
\subsubsection{科目の中でのこのコマの位置づけ}
ここまで\verb|generated quantities|ブロックの活用で，事後分布，事後予測分布を生成できること，パラメータの世界，データの世界それぞれでの仮説が検証できることを見てきた。
またパラメータリカバリの方法を見ることで，データ生成モデルを分析前に活用する方法についても見た。

続いて多群の平均値差を求めるモデルを考える。まずは一要因3水準のBetweenモデルから，3つの平均値をバラバラに求めること，生成量から差分を計算することを考える。データやパラメータレベルでの仮説的な検証方法を再確認し，帰無仮説検定の枠組みで考えなければならなかったアルファ水準のインフレ問題が生じないことを，確率の考え方の違いに即して理解する。
続いて差をモデルに組み込む方法を考える。ここでパラメータの数に制約をかける方法として\verb|transformed parameters|ブロックの使い方を導入する。
またパラメータの数の制約は，自由度の概念と深く関係していることへの洞察を得る。またパラメータリカバリのコードがリバースエンジニアリングのコードと同じもので，鏡合わせの関係にあることを確認する。
続いて交互作用が含まれるモデルを考える。ここで制約からどれだけのパラメータが必要か，どのように組み上げるかを学ぶことができる。


\subsubsection{コマ主題細目}
\begin{description}
	\item[要因計画モデル] 一要因3水準Betweenモデルを考える。対応のない二群の時のように，これは素直に3群のモデルとして表現できるし，群間の差を生成量として計算できることを再確認する。また検定と違って差の大きさを直接検証すること，確率的判断を含まないことから，アルファ水準のインフレ問題に悩む必要がないことがわかる。この考え方は，確率の捉え方の違いにもつながることに留意する。
	\item[パラメータの変形と制約] 3群のうち，ある群を基準にした差分を直接パラメータとして推定するモデルを考えると，2つの差分を計算することができる。またある群を基準におかなくとも，全体平均を基準に置くことができるが，その場合は差分のパラメータに制約をかける必要がある。これらの制約を含んだモデルを，\verb|transformed parameters|ブロックを使って作ることを確認する。ここでパラメータの自由度について理解する。
	\item[モデルの洗練] 技術的な問題であるが，多群モデルの場合は一般的に描画するためにも，整然データの形式に整えておくことが望ましい。書いたモデルの一般化という観点から，変数にできるところは変数にするなど，コードの洗練を試みる。ここで群を識別する変数を導入することは，今後の個人内反復測定モデルにも応用できる点であるので，しっかり理解する。
	\item[パラメータリカバリ] これらのモデルのパラメータリカバリから，仮想データはデータ生成モデルを逆転させるだけで出来上がことがわかり，要因計画を裏側から眺めるような，新しい観点からの理解が進むと考えられる。X
\end{description}
\subsubsection{キーワード}
\begin{itemize}
	\item 要因計画
	\item 整然データ
	\item 生成量
	\item パラメータリカバリ
\end{itemize}
\subsection{授業情報}
\paragraph{コマの展開方法}
講義/遠隔/演習
\subsubsection{予習・復習課題}
\paragraph{予習}パラメータリカバリの必要性など，データ生成モデルのアプローチにおける標準的な手順を再確認する。また対応のない二群の平均値差の検定，要因計画の検定についてこれまでの復習をしておくことで，今回の内容の理解が深まるだろう。
\paragraph{復習}要因数が増えた場合どのようになるか，またその都度Stanモデルを書き換えなくても良くなるような一般的な書き方について，自分なりに試行錯誤することが望ましい。